\section{Experiments}
\label{sec:experiments}

% ---------------------------------------------------------------------------
\subsection{Baselines}
\label{sec:baselines}

We compare against four baselines spanning the mixed-integer search
space our method addresses.
B0, B1, and B2 fix seam positions at the midpoints of their respective
landmark quads ($\boldsymbol{\delta} = 0.5$) and differ only in how they
assign texture phase offsets, covering the full range of
$\boldsymbol{\kappa}$-only strategies from no texture awareness (B0) to
the best budget-matched $\boldsymbol{\kappa}$ assignment (B2).
B3 samples the full joint search space of the GA, $(\boldsymbol{\delta},
\boldsymbol{\kappa})$, at the same evaluation budget as one GA run,
covering the "is structured search necessary" question with uniform
random sampling.
All baselines include Stage~2 global phase alignment
(Section~\ref{sec:stage2}) for a fair comparison.

\paragraph{B3: Random search at matched budget.}
B3 draws uniform random $(\boldsymbol{\delta}, \boldsymbol{\kappa})$
samples from the same distribution as the GA's initial population
($\boldsymbol{\delta}_j \sim \mathcal{U}[0.2, 0.8]$,
$\kappa_i \sim \mathcal{U}\{0, \dots, K{-}1\}$), evaluates each with the
same geometry pipeline, Stage~2 solver, and nesting engine as the GA,
and retains the best-found individual under the same fitness weighting.
The sample count is set equal to the GA's evaluation budget
($\mathrm{pop} \times \mathrm{gens} = 2{,}000$), and because both methods
pay the same $55$--$60\,\%$ infeasibility tax on attempted evaluations
(Section~\ref{sec:ea}), they also face the same effective valid-evaluation
budget. B3 and the GA differ only in how they move through the search
space, so any improvement of the GA over B3 is attributable to selection
pressure concentrating the GA's valid samples near previously-discovered
high-quality regions, which uniform random sampling cannot exploit by
construction.

\paragraph{B2: Optimal phase selection.}
Phase offsets are assigned to minimise total seam phase mismatch without
seam repositioning.
For the pants garment ($M=4$, $K^M = 8^4 = 4{,}096$), this is solved by
exhaustive enumeration.
For the onesie ($M=11$, $K^M = 8^{11} \approx 8.6 \times 10^9$),
exhaustive search is intractable; B2 instead evaluates $5 \times 10^5$
uniformly random $\boldsymbol{\kappa}$ assignments and keeps the best.
B2 serves as the ablation for the $\boldsymbol{\delta}$ component: any
improvement of the GA over B2 on $f_\text{align}$ is attributable to seam
position optimisation.

\paragraph{B1: Greedy phase selection.}
Patches are placed in breadth-first order over the seam adjacency graph.
Each patch's phase offset $\kappa_i$ is chosen greedily from $K$ candidates
to minimise mismatch with already-placed neighbours.
This captures the best texture alignment achievable by a greedy strategy
without global optimisation or seam-position search.

\paragraph{B0: Pure greedy nesting.}
All phase offsets are set to zero ($\boldsymbol{\kappa} = \mathbf{0}$),
patches are placed in descending order of area using the bottom-left heuristic,
and no grain rotation is applied.
This represents the classical industry approach: optimising fabric efficiency
with no texture awareness.

% ---------------------------------------------------------------------------
\subsection{Implementation Details}
\label{sec:implementation}

The GA uses a population of 100 individuals evolved for 20 generations with
elite count 4 and tournament size 4.
Crossover is applied with probability 0.7 and mutation with probability 0.7.
Per-variable mutation rates are: $p_\kappa = 0.35$ (kappa flip),
$p_\rho = 0.15$ (rho flip).
Delta is mutated by Gaussian perturbation with $\sigma_\delta = 0.20$,
clamped to the range $[0.2, 0.8]$ to keep keypoints away from degenerate
mesh regions while still allowing meaningful seam repositioning.
We use $K=8$ phase bins per period.

The fabric roll width is $W = 1{,}500\,\text{mm}$ with a texture period of
$50 \times 50\,\text{mm}$.
Fitness weights are $w_\text{fabric} = 1$, $w_\text{align} = 2.5$, $w_\text{fit} = 10$.
All experiments run $N \in \{5, 10, 25, 50, 100\}$ bodies per roll and are
repeated with 10 independent random seeds; we report mean $\pm$ standard
deviation.
Each fitness evaluation takes approximately 5\,s on a single CPU core
(geometry: ${\sim}$4.5\,s, Stage~2: ${\sim}$0.1\,s, nesting: $<$0.01\,s),
so a full GA run of 2{,}000 evaluations completes in approximately
3 hours.
B2 on the onesie runs its geometry pipeline once and then scores
$5 \times 10^5$ $\boldsymbol{\kappa}$ assignments at fixed
$\boldsymbol{\delta} = 0.5$, inspecting 250$\times$ more configurations
than the GA in a strictly smaller subspace; the GA nonetheless
dominates B2 on $f_\text{align}$ (Section~\ref{sec:results}),
justifying the cost of population-based search over the full joint
space.

% ---------------------------------------------------------------------------
\subsection{Garment Types}
\label{sec:garments}

We evaluate on three garment types spanning a range of topological
complexity (Table~\ref{tab:results_garments}).

\textbf{Onesie with sleeves} (\texttt{onesie\_sleeves}) is our primary
evaluation garment: $M=11$ patches connected by 26 geodesic seams, of
which 16 carry non-zero importance weights.
The combination of torso, legs, and sleeves produces a rich seam network
where stripe alignment across front/back panels, waist transitions,
and sleeve attachments must all be resolved jointly.
All ablations and detailed body-count sweeps
(Table~\ref{tab:results_onesie}) are reported on this garment.

\textbf{Pants} (\texttt{lower}) consists of $M=4$ patches connected by
6 active geodesic seams.
The simpler topology ($K^M = 4{,}096$) allows exhaustive B2 search.
On this garment, B0, B1, and B2 produce identical $f_\text{align}$
because Stage~2 fully resolves all $\boldsymbol{\kappa}$ differences
at $M=4$; the GA's improvement is therefore attributable entirely to
$\boldsymbol{\delta}$ search.

\textbf{Sleeveless shirt} (\texttt{upper}) is the simplest garment
with $M=2$ patches and a single active seam.
It serves as a lower-complexity sanity check: the GA should improve
alignment even when the seam network is minimal.

% ---------------------------------------------------------------------------
\subsection{Results}
\label{sec:results}

Table~\ref{tab:results_onesie} reports detailed results for the onesie
garment (primary evaluation) and Table~\ref{tab:results_garments}
summarises the GA's improvement across all three garment types.
The GA reduces $f_\text{align}$ relative to every baseline at every
body count, at a small cost in $f_\text{fabric}$ of a few percent
relative to the strongest baselines, as encoded by the fitness weights
$w_\text{fabric} = 1$, $w_\text{align} = 2.5$.
The fit guard $f_\text{fit}$ is zero on the majority of runs and remains
a small fraction of baseline patch area in the remainder$^\dagger$, so the
reported $f_\text{align}$ gains cannot be attributed to garment shrinkage.

Three comparisons decompose the GA's $f_\text{align}$ advantage into
individually attributable components.
On the onesie at $N=100$, the progression
B2 ($f_\text{align} = 0.90$)
$\rightarrow$ B3 ($0.79$)
$\rightarrow$ $\text{GA}_{\text{no-}\kappa}$ ($0.77$)
$\rightarrow$ GA ($0.72$)
shows that each mechanism contributes measurably.

(i) \emph{Does $\boldsymbol{\delta}$ search help?}
B2 fixes $\boldsymbol{\delta} = 0.5$ and searches only
$\boldsymbol{\kappa}$; B3 searches both
$\boldsymbol{\delta}$ and $\boldsymbol{\kappa}$ uniformly at the same
evaluation budget.
The 12\% drop from B2 to B3 shows that
$\boldsymbol{\delta}$ optimisation is the single largest source of
$f_\text{align}$ improvement, even when performed by random sampling.

(ii) \emph{Do evolutionary operators help beyond random search?}
B3 and the GA pay the same infeasibility tax
(Section~\ref{sec:ea}): 55--60\,\% of attempted evaluations fail with
a death-penalty fitness in both cases, so neither method enjoys a larger
effective valid-evaluation budget.
The GA's improvement over B3 (${\sim}$8\% on $f_\text{align}$) is
therefore attributable to \emph{where} in the feasibility region its
valid samples fall: selection and mutation around elite parents
concentrate sampling near the current best, which uniform random
sampling cannot exploit.

(iii) \emph{Does $\boldsymbol{\kappa}$ mutation help on top of
$\boldsymbol{\delta}$ search?}
Removing $\boldsymbol{\kappa}$ mutation from the GA
($\text{GA}_{\text{no-}\kappa}$, Section~\ref{sec:ablation}) costs
an additional ${\sim}$6\% on $f_\text{align}$.
This is the mixed-integer contribution: there exist global phase
reconfigurations that Stage~2's local rigid corrections cannot
reach, and the $\boldsymbol{\kappa}$ operator finds them.

\noindent
Taken together, the three comparisons confirm the landscape analysis
in Section~\ref{sec:ea}: $\boldsymbol{\delta}$ search provides the
largest single improvement, evolutionary operators exploit the sparse
$\boldsymbol{\delta}$ landscape more effectively than uniform random
sampling, and $\boldsymbol{\kappa}$ mutation contributes a further
refinement that no $\boldsymbol{\delta}$-only method can access.

\paragraph{Operator ablation: the role of $\boldsymbol{\kappa}$ mutation.}
\label{sec:ablation}
To isolate the contribution of the $\boldsymbol{\kappa}$ mutation
operator on top of $\boldsymbol{\delta}$ optimisation, we evaluate a GA
variant ($\text{GA}_{\text{no-}\kappa}$) in which the per-gene
$\boldsymbol{\kappa}$ flip probability is set to zero ($p_\kappa = 0$),
so that phase offsets evolve only through the initial population and
uniform crossover. All other operators, the $\boldsymbol{\delta}$
Gaussian mutation, tournament selection with Pareto dominance, and
elitism, are left unchanged.
On the onesie at $N=100$, $\text{GA}_{\text{no-}\kappa}$ achieves
$f_\text{align} \approx 0.77$, compared to the full GA's
$f_\text{align} = 0.72 \pm 0.04$: the $\boldsymbol{\kappa}$ operator
accounts for approximately 6\% of additional $f_\text{align}$
improvement, which is larger than the 3\% gain attributable to
evolutionary exploitation of $\boldsymbol{\delta}$ over random search
(B3 $\rightarrow$ $\text{GA}_{\text{no-}\kappa}$).
This confirms that Stage~2 does \emph{not} absorb the full
$\boldsymbol{\kappa}$ signal: while Stage~2's Levenberg--Marquardt
solver handles local phase corrections via small rigid translations,
there exist global phase reconfigurations (e.g.\ shifting an entire
sleeve or leg panel by a full period fraction) that only discrete
$\boldsymbol{\kappa}$ mutation can reach.
The mixed-integer operator design is therefore not redundant with
Stage~2; it addresses a complementary part of the search space.
% TODO: tighten percentages once more GA-no-kappa seeds land

\begin{table}[t]
\centering
\caption{
  Quantitative results on the \textbf{onesie with sleeves} ($M=11$).
  $f_\text{fabric}$: normalised fabric height (lower is better).
  $f_\text{align}$: seam phase mismatch (lower is better).
  B0, B1, B2 are deterministic (single value).
  B3 and $\text{GA}_{\text{no-}\kappa}$ report best-of-pool / single
  seed ($^\dagger$).
  GA reports mean $\pm$ std over 4--5 seeds (to be extended).
}
\label{tab:results_onesie}
\setlength{\tabcolsep}{5pt}
\begin{tabular}{llcc}
\toprule
$N$ & Method & $f_\text{fabric}$ & $f_\text{align}$ \\
\midrule
\multirow{6}{*}{5}
  & B0 & 1.193 & 1.607 \\
  & B1 & 1.283 & 1.159 \\
  & B2 & 1.221 & 0.902 \\
  & B3$^\dagger$ & -- & -- \\
  & $\text{GA}_{\text{no-}\kappa}{}^\dagger$ & -- & -- \\
  & GA & $1.241 \pm 0.049$ & $\mathbf{0.706} \pm 0.060$ \\
\midrule
\multirow{6}{*}{10}
  & B0 & -- & -- \\
  & B1 & -- & -- \\
  & B2 & -- & -- \\
  & B3$^\dagger$ & -- & -- \\
  & $\text{GA}_{\text{no-}\kappa}{}^\dagger$ & -- & -- \\
  & GA$^\dagger$ & -- & -- \\
\midrule
\multirow{6}{*}{25}
  & B0 & -- & -- \\
  & B1 & -- & -- \\
  & B2 & -- & -- \\
  & B3$^\dagger$ & -- & -- \\
  & $\text{GA}_{\text{no-}\kappa}{}^\dagger$ & -- & -- \\
  & GA$^\dagger$ & -- & -- \\
\midrule
\multirow{6}{*}{50}
  & B0 & -- & -- \\
  & B1 & -- & -- \\
  & B2 & -- & -- \\
  & B3$^\dagger$ & -- & -- \\
  & $\text{GA}_{\text{no-}\kappa}{}^\dagger$ & -- & -- \\
  & GA$^\dagger$ & -- & -- \\
\midrule
\multirow{6}{*}{100}
  & B0 & 1.177 & 1.607 \\
  & B1 & 1.267 & 1.159 \\
  & B2 & 1.177 & 0.902 \\
  & B3$^\dagger$ & $1.176$ & $0.791$ \\
  & $\text{GA}_{\text{no-}\kappa}{}^\dagger$ & $1.155$ & $0.770$ \\
  & GA & $1.181 \pm 0.030$ & $\mathbf{0.723} \pm 0.040$ \\
\bottomrule
\end{tabular}
\end{table}

\begin{table}[t]
\centering
\caption{
  Cross-garment comparison at $N=200$ bodies.
  $f_\text{align}$: seam phase mismatch (lower is better).
  B0 ($\boldsymbol{\kappa}=\mathbf{0}$, no texture awareness) and B2
  (best $\boldsymbol{\kappa}$ at fixed $\boldsymbol{\delta}=0.5$) bracket
  the $\boldsymbol{\kappa}$-only baseline range.
  GA: single seed shown; onesie numbers use $N=100$ data as proxy
  (to be updated with $N=200$).
  $\Delta$: GA improvement over B2 on $f_\text{align}$.
}
\label{tab:results_garments}
\setlength{\tabcolsep}{5pt}
\begin{tabular}{lrcccr}
\toprule
Garment & $M$ & B0 & B2 & GA & $\Delta$ vs B2 \\
\midrule
Shirt (\texttt{upper})
  & 2  & --    & --    & --    & -- \\
Pants (\texttt{lower})
  & 4  & 0.164 & 0.164 & $\mathbf{0.118}$ & $-28\,\%$ \\
Onesie (\texttt{onesie\_sleeves})
  & 11 & 1.607 & 0.902 & $\mathbf{0.723}$ & $-20\,\%$ \\
\bottomrule
\end{tabular}
\end{table}

\paragraph{Qualitative comparison.}
Figure~\ref{fig:comparison} shows a side-by-side comparison of the nesting
layout produced by B0 (no texture awareness) and the GA on the pants garment
with a horizontal stripe pattern.
The GA optimises phase offsets $\boldsymbol{\kappa}$ to shift the stripe
bands within each patch so that they align across seam boundaries, reducing
the visible discontinuity at each seam.
% TODO: replace with onesie comparison from final runs for a more compelling example

\begin{figure}[t]
\centering
\includegraphics[width=\linewidth]{comparison}
\caption{
  Nesting layout for the pants garment ($M=4$, $N=1$, stripes).
  Left: B0 baseline ($\boldsymbol{\kappa}=\mathbf{0}$, no texture alignment).
  Right: GA result with optimised phase offsets.
  Stripe lines within patches show how the GA shifts texture phases to
  improve alignment across seam boundaries.
}
\label{fig:comparison}
\end{figure}

\paragraph{Convergence and landscape validation.}
Figure~\ref{fig:convergence} reports the best-so-far weighted fitness
$F_\Sigma$ per generation for each body count, and serves as empirical
validation of the landscape analysis in Section~\ref{sec:ea}.
Three predictions from that analysis are testable from the convergence
data.
First, the $\boldsymbol{\delta}$ landscape is piecewise-constant over
Voronoi plateaus, so improvements should appear as discrete jumps when a
mutation crosses a cell boundary rather than as smooth descent.
Second, consecutive improvements should be separated by extended stagnant
stretches where most mutations fall inside the current plateau.
Third, the magnitude of each jump should be non-trivial because crossing a
plateau boundary corresponds to a qualitatively different LSCM flattening,
not an infinitesimal perturbation.
All three predictions hold: across all body counts and seeds, improvements
are concentrated in roughly $3$--$7$ of $20$ generations$^\dagger$, with
total drops of $14$--$29\%$ in $F_\Sigma$, matching the step structure
predicted by the Voronoi tessellation.
This is not a weakness of the optimiser: it is a structural property of
the decision variable $\boldsymbol{\delta}$ under the mesh-vertex snap,
and the population-based search is what allows the optimiser to cross
plateau boundaries at all, since any single-point local search would halt
the first time a mutation failed to improve fitness.
% TODO: replace 3-7 and 14-29 with final 10-seed numbers once runs complete

\begin{figure}[t]
\centering
\includegraphics[width=\linewidth]{convergence}
\caption{
  Best-so-far $F_\Sigma$ vs.\ generation for the onesie garment,
  median and interquartile range over 10 seeds per body count $N$.
  Improvements are concentrated in a small number of generations and
  appear as discrete jumps separated by stagnant stretches, directly
  reflecting the Voronoi tessellation of the $\boldsymbol{\delta}$
  landscape (Section~\ref{sec:ea}): mutations that do not cross a cell
  boundary produce zero fitness change, and each jump corresponds to a
  qualitatively different LSCM flattening.
  $^\dagger$Placeholder; to be replaced with final 10-seed results.
}
\label{fig:convergence}
\end{figure}

% ---------------------------------------------------------------------------
\subsection{Qualitative Results}
\label{sec:qualitative}

Figure~\ref{fig:wallpaper_groups} shows the GA result rendered on the
sleeveless shirt for all eight supported wallpaper groups.
Each garment was optimised independently with the corresponding symmetry
group, demonstrating that the method correctly adapts seam placement and
phase offsets to the specific symmetry structure of each pattern.
Stripe-like groups (PM, PG, PMG, PGG) produce clearly directed motifs that
continue coherently across seam lines, while rotationally symmetric groups
(P4, P4M) yield more isotropic patterns that are inherently less sensitive
to seam placement but still benefit from phase alignment.

\begin{figure}[t]
\centering
\includegraphics[width=\linewidth]{wallpaper_groups}
\caption{
  GA-optimised sleeveless shirt rendered with all eight supported wallpaper
  groups.
  From left to right, top row: stripes (PM), diagonal stripes (PM),
  grid (PMM), pinwheel (P4).
  Bottom row: polka dots (P4M), herringbone (PG), chevron (PMG),
  brick bond (PGG).
  Each result uses the best individual found by the GA in a single run
  ($N=1$ body, pop=20, 5 generations).
}
\label{fig:wallpaper_groups}
\end{figure}

% ---------------------------------------------------------------------------
